% ===========================================================================
%  evnx for Mobile Developers  —  Expo / React Native, Flutter
%  Persona MK.  Build: pdflatex -interaction=nonstopmode mobile.tex
% ===========================================================================
\documentclass[aspectratio=169,10pt]{beamer}

\newcommand{\capturedir}{../captures}
% ---------------------------------------------------------------------------
%  Shared preamble for the seven audience decks.
%
%  Each deck does, before the document body starts:
%     \newcommand{\capturedir}{../captures}
%     % ---------------------------------------------------------------------------
%  Shared preamble for the seven audience decks.
%
%  Each deck does, before the document body starts:
%     \newcommand{\capturedir}{../captures}
%     % ---------------------------------------------------------------------------
%  Shared preamble for the seven audience decks.
%
%  Each deck does, before the document body starts:
%     \newcommand{\capturedir}{../captures}
%     \input{persona-preamble}
%
%  \capturedir points at ../captures so both pools resolve:
%     \capture{24-scan.txt}              — the original v0.5.0 captures
%     \capture{persona/fe-01-scan-dist.txt} — captures recorded for these decks
% ---------------------------------------------------------------------------

\input{../common/evnx-preamble}

% ---- provenance ------------------------------------------------------------
% The original decks were pinned to tag v0.5.0 (85a52c3). Everything captured
% for the persona decks was re-run against 0.5.2, so the footline differs.
\renewcommand{\verifiedon}{\tiny\textcolor{evnxMuted}{verified against evnx 0.5.2 --- real run, not transcribed}}

% ---- audience badge on the title slide ------------------------------------
\newcommand{\audience}[1]{%
  \begin{center}\small\textcolor{evnxAccent}{\textbf{Audience: #1}}\end{center}}

% ---- verdict chips ---------------------------------------------------------
\newcommand{\solved}{\textcolor{evnxOk}{\ensuremath{\checkmark}}}
\newcommand{\partialmark}{\textcolor{evnxWarn}{\textbf{\textasciitilde}}}
\newcommand{\unsolved}{\textcolor{evnxErr}{\ensuremath{\times}}}

% ---- a correction block: something the source material got wrong -----------
\newenvironment{correction}[1]{%
  \begin{alertblock}{\small Correction --- #1}\small}{\end{alertblock}}

% ---- a trap block: works in the demo, fails in real use --------------------
\newenvironment{trap}[1]{%
  \setbeamercolor{block title}{bg=evnxWarn,fg=white}%
  \setbeamercolor{block body}{bg=evnxWarn!8}%
  \begin{block}{\small #1}\small}{\end{block}}

% ---- the persona card ------------------------------------------------------
\newcommand{\personacard}[4]{%
  \begin{block}{#1}
    \small
    \textbf{Stack:} #2 \\
    \textbf{Ships:} #3 \\
    \textbf{Loses time to:} #4
  \end{block}}

% ---- speaker note (rendered small and muted, not hidden) -------------------
\newcommand{\saythis}[1]{%
  \vspace{0.4em}{\footnotesize\textcolor{evnxMuted}{\textit{Say:} #1}}}

%
%  \capturedir points at ../captures so both pools resolve:
%     \capture{24-scan.txt}              — the original v0.5.0 captures
%     \capture{persona/fe-01-scan-dist.txt} — captures recorded for these decks
% ---------------------------------------------------------------------------

% ---------------------------------------------------------------------------
%  evnx presentation preamble — shared by all three decks
%
%  Engine: pdfLaTeX (also builds under XeLaTeX / LuaLaTeX).
%  Packages: beamer, listings, xcolor, amssymb, newunicodechar, booktabs,
%            fancyvrb, hyperref — all in texlive-latex-recommended + beamer.
%
%  \capturedir must be defined by the including document BEFORE \input of
%  this file, e.g.  \newcommand{\capturedir}{../captures}
% ---------------------------------------------------------------------------

\usetheme{Madrid}
\usecolortheme{seahorse}
\usefonttheme{professionalfonts}
\setbeamertemplate{navigation symbols}{}
\setbeamertemplate{caption}[numbered]
\setbeamerfont{frametitle}{size=\large}
\setbeamerfont{title}{size=\LARGE,series=\bfseries}

\usepackage[T1]{fontenc}
\usepackage[utf8]{inputenc}
\usepackage{lmodern}
\usepackage{amssymb}
\usepackage{booktabs}
\usepackage{xcolor}
\usepackage{listings}
\usepackage{fancyvrb}
\usepackage{newunicodechar}
\usepackage{tikz}
\usetikzlibrary{arrows.meta,positioning,shapes.geometric,fit,backgrounds}

% ---- palette ---------------------------------------------------------------
\definecolor{evnxInk}{HTML}{1F2933}
\definecolor{evnxAccent}{HTML}{2F6F4E}
\definecolor{evnxWarn}{HTML}{B7791F}
\definecolor{evnxErr}{HTML}{B03A2E}
\definecolor{evnxOk}{HTML}{2F6F4E}
\definecolor{evnxMuted}{HTML}{6B7280}
\definecolor{evnxTermBg}{HTML}{F5F5F2}
\definecolor{evnxRule}{HTML}{D1D5DB}

\setbeamercolor{structure}{fg=evnxAccent}
\setbeamercolor{alerted text}{fg=evnxErr}

% ---- glyph mapping for body text ------------------------------------------
% evnx emits exactly these 11 non-ASCII characters (enumerated from the
% captured logs, not guessed).
\newcommand{\evnxCheck}{\textcolor{evnxOk}{\ensuremath{\checkmark}}}
\newcommand{\evnxCross}{\textcolor{evnxErr}{\ensuremath{\times}}}
\newcommand{\evnxWarnSign}{\textcolor{evnxWarn}{\textbf{!}}}

\newunicodechar{✓}{\evnxCheck}
\newunicodechar{✗}{\evnxCross}
\newunicodechar{⚠}{\evnxWarnSign}
\newunicodechar{→}{\ensuremath{\rightarrow}}
\newunicodechar{↔}{\ensuremath{\leftrightarrow}}
\newunicodechar{─}{\textendash}
\newunicodechar{—}{---}
\newunicodechar{·}{\textperiodcentered}
\newunicodechar{…}{\ldots}
\newunicodechar{•}{\textbullet}
\newunicodechar{️}{}                % U+FE0F variation selector — invisible

% ---- terminal listing style ------------------------------------------------
% `literate` is required: newunicodechar does not apply inside listings.
\lstdefinestyle{evnxterm}{
  basicstyle=\ttfamily\scriptsize,
  backgroundcolor=\color{evnxTermBg},
  frame=single,
  rulecolor=\color{evnxRule},
  framesep=4pt,
  breaklines=true,
  breakatwhitespace=false,
  columns=fullflexible,
  keepspaces=true,
  showstringspaces=false,
  upquote=true,
  extendedchars=true,
  literate=
    {✓}{{\textcolor{evnxOk}{$\checkmark$}}}1
    {✗}{{\textcolor{evnxErr}{$\times$}}}1
    {⚠}{{\textcolor{evnxWarn}{\textbf{!}}}}1
    {→}{{$\rightarrow$}}1
    {↔}{{$\leftrightarrow$}}1
    {─}{{-}}1
    {—}{{---}}1
    {·}{{$\cdot$}}1
    {…}{{\ldots}}1
    {•}{{\textbullet}}1
    {️}{{}}0
    {é}{{\'e}}1
}

\lstdefinestyle{evnxcode}{
  style=evnxterm,
  basicstyle=\ttfamily\scriptsize,
  backgroundcolor=\color{white},
}

\lstset{style=evnxterm}

% ---- helper macros ---------------------------------------------------------
% \capture{file}        — include a real captured run verbatim
% \capturepart{f}{a}{b} — include lines a..b of a captured run
\newcommand{\capture}[1]{\lstinputlisting[style=evnxterm]{\capturedir/#1}}
\newcommand{\capturepart}[3]{%
  \lstinputlisting[style=evnxterm,firstline=#2,lastline=#3]{\capturedir/#1}}

% a labelled "real output" box
\newcommand{\realrun}[1]{%
  \begin{center}\scriptsize\textcolor{evnxMuted}{real run \textendash\ #1}\end{center}}

\newcommand{\cmd}[1]{\texttt{\textbf{#1}}}
\newcommand{\flag}[1]{\texttt{#1}}
\newcommand{\file}[1]{\texttt{#1}}

% exit-code chips
\newcommand{\exitzero}{\textcolor{evnxOk}{\texttt{0}}}
\newcommand{\exitone}{\textcolor{evnxWarn}{\texttt{1}}}
\newcommand{\exittwo}{\textcolor{evnxErr}{\texttt{2}}}

% a standard "verified against" footline for factual slides
\newcommand{\verifiedon}{\tiny\textcolor{evnxMuted}{verified against evnx 0.5.0 @ 9cfe75b}}

\usepackage[colorlinks=true,linkcolor=evnxAccent,urlcolor=evnxAccent]{hyperref}


% ---- provenance ------------------------------------------------------------
% The original decks were pinned to tag v0.5.0 (85a52c3). Everything captured
% for the persona decks was re-run against 0.5.2, so the footline differs.
\renewcommand{\verifiedon}{\tiny\textcolor{evnxMuted}{verified against evnx 0.5.2 --- real run, not transcribed}}

% ---- audience badge on the title slide ------------------------------------
\newcommand{\audience}[1]{%
  \begin{center}\small\textcolor{evnxAccent}{\textbf{Audience: #1}}\end{center}}

% ---- verdict chips ---------------------------------------------------------
\newcommand{\solved}{\textcolor{evnxOk}{\ensuremath{\checkmark}}}
\newcommand{\partialmark}{\textcolor{evnxWarn}{\textbf{\textasciitilde}}}
\newcommand{\unsolved}{\textcolor{evnxErr}{\ensuremath{\times}}}

% ---- a correction block: something the source material got wrong -----------
\newenvironment{correction}[1]{%
  \begin{alertblock}{\small Correction --- #1}\small}{\end{alertblock}}

% ---- a trap block: works in the demo, fails in real use --------------------
\newenvironment{trap}[1]{%
  \setbeamercolor{block title}{bg=evnxWarn,fg=white}%
  \setbeamercolor{block body}{bg=evnxWarn!8}%
  \begin{block}{\small #1}\small}{\end{block}}

% ---- the persona card ------------------------------------------------------
\newcommand{\personacard}[4]{%
  \begin{block}{#1}
    \small
    \textbf{Stack:} #2 \\
    \textbf{Ships:} #3 \\
    \textbf{Loses time to:} #4
  \end{block}}

% ---- speaker note (rendered small and muted, not hidden) -------------------
\newcommand{\saythis}[1]{%
  \vspace{0.4em}{\footnotesize\textcolor{evnxMuted}{\textit{Say:} #1}}}

%
%  \capturedir points at ../captures so both pools resolve:
%     \capture{24-scan.txt}              — the original v0.5.0 captures
%     \capture{persona/fe-01-scan-dist.txt} — captures recorded for these decks
% ---------------------------------------------------------------------------

% ---------------------------------------------------------------------------
%  evnx presentation preamble — shared by all three decks
%
%  Engine: pdfLaTeX (also builds under XeLaTeX / LuaLaTeX).
%  Packages: beamer, listings, xcolor, amssymb, newunicodechar, booktabs,
%            fancyvrb, hyperref — all in texlive-latex-recommended + beamer.
%
%  \capturedir must be defined by the including document BEFORE \input of
%  this file, e.g.  \newcommand{\capturedir}{../captures}
% ---------------------------------------------------------------------------

\usetheme{Madrid}
\usecolortheme{seahorse}
\usefonttheme{professionalfonts}
\setbeamertemplate{navigation symbols}{}
\setbeamertemplate{caption}[numbered]
\setbeamerfont{frametitle}{size=\large}
\setbeamerfont{title}{size=\LARGE,series=\bfseries}

\usepackage[T1]{fontenc}
\usepackage[utf8]{inputenc}
\usepackage{lmodern}
\usepackage{amssymb}
\usepackage{booktabs}
\usepackage{xcolor}
\usepackage{listings}
\usepackage{fancyvrb}
\usepackage{newunicodechar}
\usepackage{tikz}
\usetikzlibrary{arrows.meta,positioning,shapes.geometric,fit,backgrounds}

% ---- palette ---------------------------------------------------------------
\definecolor{evnxInk}{HTML}{1F2933}
\definecolor{evnxAccent}{HTML}{2F6F4E}
\definecolor{evnxWarn}{HTML}{B7791F}
\definecolor{evnxErr}{HTML}{B03A2E}
\definecolor{evnxOk}{HTML}{2F6F4E}
\definecolor{evnxMuted}{HTML}{6B7280}
\definecolor{evnxTermBg}{HTML}{F5F5F2}
\definecolor{evnxRule}{HTML}{D1D5DB}

\setbeamercolor{structure}{fg=evnxAccent}
\setbeamercolor{alerted text}{fg=evnxErr}

% ---- glyph mapping for body text ------------------------------------------
% evnx emits exactly these 11 non-ASCII characters (enumerated from the
% captured logs, not guessed).
\newcommand{\evnxCheck}{\textcolor{evnxOk}{\ensuremath{\checkmark}}}
\newcommand{\evnxCross}{\textcolor{evnxErr}{\ensuremath{\times}}}
\newcommand{\evnxWarnSign}{\textcolor{evnxWarn}{\textbf{!}}}

\newunicodechar{✓}{\evnxCheck}
\newunicodechar{✗}{\evnxCross}
\newunicodechar{⚠}{\evnxWarnSign}
\newunicodechar{→}{\ensuremath{\rightarrow}}
\newunicodechar{↔}{\ensuremath{\leftrightarrow}}
\newunicodechar{─}{\textendash}
\newunicodechar{—}{---}
\newunicodechar{·}{\textperiodcentered}
\newunicodechar{…}{\ldots}
\newunicodechar{•}{\textbullet}
\newunicodechar{️}{}                % U+FE0F variation selector — invisible

% ---- terminal listing style ------------------------------------------------
% `literate` is required: newunicodechar does not apply inside listings.
\lstdefinestyle{evnxterm}{
  basicstyle=\ttfamily\scriptsize,
  backgroundcolor=\color{evnxTermBg},
  frame=single,
  rulecolor=\color{evnxRule},
  framesep=4pt,
  breaklines=true,
  breakatwhitespace=false,
  columns=fullflexible,
  keepspaces=true,
  showstringspaces=false,
  upquote=true,
  extendedchars=true,
  literate=
    {✓}{{\textcolor{evnxOk}{$\checkmark$}}}1
    {✗}{{\textcolor{evnxErr}{$\times$}}}1
    {⚠}{{\textcolor{evnxWarn}{\textbf{!}}}}1
    {→}{{$\rightarrow$}}1
    {↔}{{$\leftrightarrow$}}1
    {─}{{-}}1
    {—}{{---}}1
    {·}{{$\cdot$}}1
    {…}{{\ldots}}1
    {•}{{\textbullet}}1
    {️}{{}}0
    {é}{{\'e}}1
}

\lstdefinestyle{evnxcode}{
  style=evnxterm,
  basicstyle=\ttfamily\scriptsize,
  backgroundcolor=\color{white},
}

\lstset{style=evnxterm}

% ---- helper macros ---------------------------------------------------------
% \capture{file}        — include a real captured run verbatim
% \capturepart{f}{a}{b} — include lines a..b of a captured run
\newcommand{\capture}[1]{\lstinputlisting[style=evnxterm]{\capturedir/#1}}
\newcommand{\capturepart}[3]{%
  \lstinputlisting[style=evnxterm,firstline=#2,lastline=#3]{\capturedir/#1}}

% a labelled "real output" box
\newcommand{\realrun}[1]{%
  \begin{center}\scriptsize\textcolor{evnxMuted}{real run \textendash\ #1}\end{center}}

\newcommand{\cmd}[1]{\texttt{\textbf{#1}}}
\newcommand{\flag}[1]{\texttt{#1}}
\newcommand{\file}[1]{\texttt{#1}}

% exit-code chips
\newcommand{\exitzero}{\textcolor{evnxOk}{\texttt{0}}}
\newcommand{\exitone}{\textcolor{evnxWarn}{\texttt{1}}}
\newcommand{\exittwo}{\textcolor{evnxErr}{\texttt{2}}}

% a standard "verified against" footline for factual slides
\newcommand{\verifiedon}{\tiny\textcolor{evnxMuted}{verified against evnx 0.5.0 @ 9cfe75b}}

\usepackage[colorlinks=true,linkcolor=evnxAccent,urlcolor=evnxAccent]{hyperref}


% ---- provenance ------------------------------------------------------------
% The original decks were pinned to tag v0.5.0 (85a52c3). Everything captured
% for the persona decks was re-run against 0.5.2, so the footline differs.
\renewcommand{\verifiedon}{\tiny\textcolor{evnxMuted}{verified against evnx 0.5.2 --- real run, not transcribed}}

% ---- audience badge on the title slide ------------------------------------
\newcommand{\audience}[1]{%
  \begin{center}\small\textcolor{evnxAccent}{\textbf{Audience: #1}}\end{center}}

% ---- verdict chips ---------------------------------------------------------
\newcommand{\solved}{\textcolor{evnxOk}{\ensuremath{\checkmark}}}
\newcommand{\partialmark}{\textcolor{evnxWarn}{\textbf{\textasciitilde}}}
\newcommand{\unsolved}{\textcolor{evnxErr}{\ensuremath{\times}}}

% ---- a correction block: something the source material got wrong -----------
\newenvironment{correction}[1]{%
  \begin{alertblock}{\small Correction --- #1}\small}{\end{alertblock}}

% ---- a trap block: works in the demo, fails in real use --------------------
\newenvironment{trap}[1]{%
  \setbeamercolor{block title}{bg=evnxWarn,fg=white}%
  \setbeamercolor{block body}{bg=evnxWarn!8}%
  \begin{block}{\small #1}\small}{\end{block}}

% ---- the persona card ------------------------------------------------------
\newcommand{\personacard}[4]{%
  \begin{block}{#1}
    \small
    \textbf{Stack:} #2 \\
    \textbf{Ships:} #3 \\
    \textbf{Loses time to:} #4
  \end{block}}

% ---- speaker note (rendered small and muted, not hidden) -------------------
\newcommand{\saythis}[1]{%
  \vspace{0.4em}{\footnotesize\textcolor{evnxMuted}{\textit{Say:} #1}}}


\title{evnx for mobile developers}
\subtitle{Expo and Flutter --- where ``the binary is public'' changes everything}
\date{}

\begin{document}

\begin{frame}[plain]
  \titlepage
  \audience{Expo / React Native · Flutter flavours · EAS · Fastlane}
  \begin{center}\footnotesize
  \textcolor{evnxMuted}{evnx 0.5.2. Every terminal block is a real run.}
  \end{center}
\end{frame}

% =========================================================================
\begin{frame}{This deck in one slide}
  \begin{block}{The rule your team should adopt, whatever tool you pick}
    \textbf{The vault for an app build should contain only values that are
    safe to be public.} Real secrets stay on the backend.
  \end{block}

  \vspace{0.7em}
  \begin{block}{Because everything else follows from it}
    An APK is a zip file. An IPA is a zip file. Anything your build reads at
    compile time is extractable by anyone who downloads your app. No secret
    manager changes that --- it only changes who can get a secret
    \emph{before} the build.
  \end{block}

  \saythis{If the room accepts this slide, the rest is logistics. If they do
  not, nothing later will help.}
\end{frame}

\begin{frame}{Contents}\tableofcontents[hideallsubsections]\end{frame}

% =========================================================================
\section{Three releases that went wrong}

\begin{frame}{You have shipped these}
  \small
  \begin{block}{The asset}
    \texttt{flutter\_dotenv} loads \file{.env} as an \textbf{app asset}. A
    security review unpacks the APK and finds it --- including a third-party
    API secret meant only for the backend.
  \end{block}

  \begin{block}{The wrong profile}
    A staging build points at the production API.
    \texttt{EXPO\_PUBLIC\_API\_URL} was set in EAS for ``production'' and the
    build used the ``preview'' profile.
  \end{block}

  \begin{block}{The keystore}
    The Android keystore password lives in a teammate's
    \file{\textasciitilde/.gradle/gradle.properties}. They are on vacation.
    The release is blocked.
  \end{block}
\end{frame}

\begin{frame}{Which of the three evnx actually fixes}
  \small
  \begin{tabular}{@{}llp{5.8cm}@{}}
    \toprule
    \textbf{Incident} & \textbf{Verdict} & \textbf{Why} \\
    \midrule
    Bundled \file{.env} asset & \partialmark & \cmd{scan assets/} finds it ---
      \emph{after} you have already built it. No prefix awareness, no APK unpacking \\
    Wrong profile              & \partialmark & \cmd{validate -{}-env-name} per
      file; cross-file consistency is a check you write \\
    Keystore password          & \solved      & vault the \textbf{password};
      the \file{.jks} file itself is not storable \\
    \bottomrule
  \end{tabular}

  \vspace{0.9em}
  \begin{block}{Honest framing}
    The keystore one is a clean win and it unblocks releases. The other two
    evnx helps with but does not close --- and it is worth saying which is
    which before someone finds out on their own.
  \end{block}
\end{frame}

% =========================================================================
\section{The good part: one vault per flavour}

\begin{frame}[fragile]{Flavours map onto vaults directly}
\begin{lstlisting}[style=evnxcode]
evnx vault create field-app --env dev
evnx vault create field-app --env staging
evnx vault create field-app --env prod

# Flutter
evnx cloud run -- flutter run --flavor dev

# Expo
evnx cloud run -- npx expo start
evnx cloud run -- eas build --local

# Fastlane -- and this is where the keystore password comes from
evnx cloud run -- bundle exec fastlane android release
\end{lstlisting}

  \vspace{0.5em}
  {\footnotesize \cmd{cloud run} means no \file{.env} on the build machine and
  no \texttt{gradle.properties} on a colleague's laptop as the single point of
  failure.}
\end{frame}

\begin{frame}[fragile]{Flutter's \flag{-{}-dart-define-from-file} fits perfectly}
  \capture{persona/mob-02-dartdefine.txt}
  \realrun{\texttt{evnx convert -{}-to json -{}-exclude "*\_SECRET*"}}

  \vspace{0.4em}
\begin{lstlisting}[style=evnxcode]
evnx convert --to json --exclude "*_SECRET*" > config/dev.json
flutter run --flavor dev --dart-define-from-file=config/dev.json
\end{lstlisting}

  \vspace{0.3em}
  \saythis{``Look carefully at that output before the next slide.''}
\end{frame}

\begin{frame}[fragile]{Now look at what came through}
\begin{lstlisting}[style=evnxcode]
{
  "EXPO_PUBLIC_API_URL":       "https://api.acme.io",
  "EXPO_PUBLIC_MAPS_KEY":      "AIzaSyD-...",
  "ANDROID_KEYSTORE_PASSWORD": "hunter2correcthorse"   <-- into the build
}
\end{lstlisting}

  \vspace{0.6em}
  \begin{trap}{\flag{-{}-exclude} is a denylist, and denylists leak}
    \texttt{ANDROID\_KEYSTORE\_PASSWORD} does not match \texttt{*\_SECRET*}, so
    it went straight into the JSON that gets compiled into the app. The
    recommended command is the one that leaked it.
  \end{trap}

  \vspace{0.5em}
  \begin{block}{Use an allowlist instead}
    \begin{center}\ttfamily\small
    evnx convert -{}-to json -{}-include "EXPO\_PUBLIC\_*" > config/dev.json
    \end{center}
    Name what may be public. Never enumerate what may not.
  \end{block}
\end{frame}

% =========================================================================
\section{The public-surface problem}

\begin{frame}[fragile]{\cmd{validate} has no idea what \texttt{EXPO\_PUBLIC\_} means}
  \capture{persona/mob-03-validate.txt}
  \realrun{a \file{.env} holding a live Stripe key and a keystore password}

  \vspace{0.4em}
  {\footnotesize Exit \exitzero. \cmd{validate} checks file hygiene --- missing
  variables, placeholders, weak secrets, boolean traps. ``This name has a
  public prefix and this value looks secret'' is a rule that does not exist
  yet.}
\end{frame}

\begin{frame}[fragile]{\cmd{scan} is the one that helps}
  \capturepart{persona/mob-01-scan-assets.txt}{1}{20}
  \realrun{\texttt{evnx scan assets/} --- the Flutter bundled \file{.env}}

  \vspace{0.3em}
  {\footnotesize Three findings by three different mechanisms: a \textbf{value}
  pattern (Stripe, high), a \textbf{variable name} (keystore password, high),
  and \textbf{entropy} (the Maps key, low). Note the annotation on the middle
  one --- \emph{``flagged by variable name --- the value was not checked''}.}
\end{frame}

\begin{frame}{The Firebase caveat, before someone raises it}
  \begin{block}{\file{google-services.json} and \file{GoogleService-Info.plist}}
    evnx may flag keys in these by pattern. Firebase API keys are
    \textbf{meant} to be embedded in clients --- you restrict them in the
    console, you do not hide them. evnx cannot know that, so expect noise and
    suppress it deliberately rather than by turning the check off.
  \end{block}

  \vspace{0.9em}
  \begin{block}{What this tells you about the tool}
    Name-based and entropy-based detection cannot distinguish ``public by
    design'' from ``leaked''. That judgement stays with you. A tool that
    claimed otherwise would be lying.
  \end{block}
\end{frame}

\begin{frame}[fragile]{And one place \cmd{scan} will not look}
  \capture{persona/fe-01-scan-dist.txt}
  \realrun{\file{dist/bundle.js} contains a live key --- \texttt{expo export} writes here}

  \vspace{0.4em}
  \begin{trap}{\file{dist/} is on the default exclusion list}
    \cmd{npx expo export \&\& evnx scan dist/} reports success over ground it
    never looked at. Copy the output somewhere else first:
    \vspace{0.3em}
    \begin{center}\ttfamily\small
    npx expo export \&\& cp -r dist /tmp/scanme \&\& evnx scan /tmp/scanme
    \end{center}
    \vspace{0.2em}
    Tracked \textbf{NEW-1, P0}. \file{assets/} is not excluded, which is why
    the Flutter case on the previous slide worked.
  \end{trap}
\end{frame}

% =========================================================================
\section{File secrets and platforms}

\begin{frame}{What a vault cannot hold}
  \small
  \begin{tabular}{@{}lll@{}}
    \toprule
    \textbf{Artefact} & \textbf{In a vault?} & \textbf{Today} \\
    \midrule
    Keystore \emph{password}       & \solved   & a normal variable \\
    \file{.jks} / \file{.p12} file & \unsolved & keep in CI secure files \\
    App Store Connect \file{.p8}   & \unsolved & same \\
    \file{google-services.json}    & \partialmark & inline as a single-quoted multi-line value, or keep as a file \\
    \bottomrule
  \end{tabular}

  \vspace{0.9em}
  \begin{block}{The multi-line detail, since it is commonly misreported}
    evnx \textbf{does} parse multi-line values. A \textbf{single}-quoted block
    holding JSON round-trips correctly. A \textbf{double}-quoted block
    containing escaped \texttt{\textbackslash "} fails to parse --- so wrap
    JSON in single quotes.
  \end{block}
\end{frame}

\begin{frame}{EAS, Codemagic, Bitrise}
  \begin{center}\large No targets. This is a real gap.\end{center}

  \vspace{1em}
  \begin{block}{The bridge that works today}
\begin{center}\ttfamily\small
evnx convert -{}-to shell -{}-include "EXPO\_PUBLIC\_*" \textbar{} \textbackslash \\
\ \ while read -r line; do eas env:create ...; done
\end{center}
  \end{block}

  \vspace{0.8em}
  \begin{block}{And the reason it cannot be automatic}
    A zero-knowledge server holds only ciphertext, so \emph{it} can never push
    your values to EAS. Only a machine holding your key can. Any sync is a
    script you run, not a service that runs for you --- which is the price of
    the property, not an oversight.
  \end{block}
\end{frame}

\begin{frame}[fragile]{The cross-profile check nobody ships}
\begin{lstlisting}[style=evnxcode]
# "the prod profile must not point at a staging host"
for e in dev staging prod; do
  url=$(evnx convert --env-name "$e" --to json | jq -r .EXPO_PUBLIC_API_URL)
  echo "$e -> $url"
done
# then assert prod matches ^https://api\.acme\.io
\end{lstlisting}

  \vspace{0.6em}
  \begin{block}{Why this is worth ten minutes}
    ``Staging build hit production'' is the incident from slide 5. evnx has no
    cross-file consistency rules --- it validates one file at a time --- so
    this loop is the entire guard. Put it in the release pipeline.
  \end{block}
\end{frame}

% =========================================================================
\section{Adoption}

\begin{frame}[fragile]{Copy-paste}
\begin{lstlisting}[style=evnxcode]
# one vault per flavour
evnx vault create field-app --env dev

# Flutter -- allowlist, not denylist
evnx convert --to json --include "EXPO_PUBLIC_*" > config/dev.json
flutter run --flavor dev --dart-define-from-file=config/dev.json

# Expo
evnx cloud run -- npx expo start
npx expo export && cp -r dist /tmp/scanme && evnx scan /tmp/scanme

# release
evnx cloud run -- bundle exec fastlane android release

# every build, both platforms
evnx scan assets/ .        # assets/ is not excluded; dist/ and build/ are
\end{lstlisting}
\end{frame}

\begin{frame}{The scorecard, honestly}
  \begin{center}
  \begin{tabular}{@{}clp{6.2cm}@{}}
    \toprule
    & \textbf{Count} & \textbf{Items} \\
    \midrule
    \solved     & 5 & vault per flavour, dart-define JSON, Expo build-time env, \file{.env.example} drift, Fastlane credentials \\
    \partialmark & 6 & bundled-asset leaks, Firebase config noise, rotation after leak, CI credentials, cross-profile consistency, keystore \emph{passwords} only \\
    \unsolved   & 2 & EAS targets, Codemagic / Bitrise targets \\
    \bottomrule
  \end{tabular}
  \end{center}

  \vspace{0.7em}
  \begin{block}{One sentence}
    evnx is a good source of truth per flavour, and its JSON output fits
    Flutter natively. Mobile's core truth --- anything in the app is public ---
    needs prefix awareness, file secrets and platform targets it does not
    have.
  \end{block}
\end{frame}

\begin{frame}{What is coming, in priority order}
  \small
  \begin{tabular}{@{}clp{5.6cm}@{}}
    \toprule
    & \textbf{Feature} & \textbf{Fixes the slide on} \\
    \midrule
    P0 & Scan an explicitly-named excluded path & \texttt{expo export} → \file{dist/} \\
    1  & Public-surface profiles (\texttt{EXPO\_PUBLIC\_}, Flutter asset \file{.env}, dart-define) & the bundled asset \\
    2  & File secrets (keystore, \file{.p8}, \file{.p12}, service JSON) materialised by \cmd{cloud run} & the keystore file \\
    3  & Artefact scanning --- unpack APK / AAB / IPA & the security review \\
    4  & EAS / Codemagic / Bitrise targets & the manual bridge \\
    5  & Cross-file consistency rules & the wrong profile \\
    \bottomrule
  \end{tabular}
\end{frame}

\begin{frame}[plain]
  \vfill
  \begin{center}
    {\LARGE\textbf{Questions}}\\[1.2em]
    {\small \url{https://www.evnx.dev/guides} · \url{https://github.com/urwithajit9/evnx}}\\[0.8em]
    {\footnotesize\textcolor{evnxMuted}{evnx 0.5.2 · every terminal block is a real run}}
  \end{center}
  \vfill
\end{frame}

\end{document}
